\documentclass[10pt,a4paper]{amsart}
\usepackage[margin=19mm]{geometry}
\usepackage{amsmath,amssymb,mathtools}
\usepackage[scr=rsfs,cal=euler]{mathalfa}
\usepackage{xfrac}
\usepackage[unicode,hidelinks,bookmarks=false]{hyperref}
\newcommand{\ie}{i.e.\ }
\newcommand{\cf}{cf.\ }
\newcommand{\resp}{respectively\ }
\newcommand{\Subsection}{\subsection}
\providecommand{\nobreakdash}{\nobreak\hbox{-}}
\setlength{\emergencystretch}{2em}
\multlinegap0pt
\usepackage{bm}
\usepackage{bbm}
\usepackage{dsfont}
\usepackage{xspace}
\usepackage{cancel}
\usepackage{mathtools}
\usepackage{amsthm}
\makeatletter
\@ifundefined{proposition}{\newtheorem{proposition}{Proposition}}{}
\makeatother
\makeatletter
\expandafter\ifx\csname algorithmblock\endcsname\relax
\newenvironment{algorithmblock}[1]{%
  \refstepcounter{algctr}\begin{figure}[t]\small\hrule\vspace{3pt}%
  \noindent\textbf{Algorithm \thealgctr: #1}\par\vspace{2pt}%
}{\vspace{2pt}\hrule\end{figure}}
\fi
\makeatother
\title[Stochastic inflows and seasonal water value at Manantali]{Stochastic inflows and seasonal water value at Manantali}
\author{Steeven B. Affognon, Babacar M. Ndiaye, Pierre Mendy, Cheikh M. F. Kebe, Sandrine Charousset, Omar Goudiaby, Thomas Ouillon}
\date{Source: arXiv:2609.27744v1 (CC BY 4.0)}
\begin{document}
\maketitle

\noindent\emph{Source and licence.} Stochastic inflows and seasonal water value at Manantali. Version arXiv:2609.27744v1 (CC BY 4.0), distributed under the Creative Commons Attribution 4.0 International licence, as recorded in the arXiv licence metadata for this exact version and shown on the abstract page https://arxiv.org/abs/2609.27744v1. Full rights notice: licenses/arxiv-2609.27744-CC-BY-4.0.txt.\par\smallskip

\noindent\textit{Editorial scope.} Counted unit: the Proposition whose source label is \texttt{prop:app-dual} in the article (source0.tex lines 1142--1166, source label \texttt{prop:app-dual}), delivered together with the article's own complete original proof of it (source0.tex lines 1168--1190, one \texttt{proof} environment of 23 lines). No mathematics is written, repaired or re-derived: statement and proof are the article's own text line for line. Required context retained uncounted: eq:app-parametric-lp (lines 1133--1138). Every definition, display and label of those lines is retained unchanged. Omitted from this package and not redistributed: 1--1132, 1139--1141, 1167--1167, 1191--1400. Nothing inside a retained excerpt is omitted or altered: every retained body is delivered exactly as the article's lines are written, so no omission receipt of any kind accompanies this package. Citations: the retained excerpts carry no citation command at all, so no bibliography entry of the article is transcribed and no reference is redistributed. Reference adaptation: none was needed and none was made; every reference inside the delivered unit resolves inside it, so no unresolved reference or citation is produced. Printed slips kept literal: no printed word, symbol, hypothesis or number of the counted statement or of its proof was altered. Layout and class: the article's own class is \texttt{article} with packages including inputenc, fontenc, lmodern, geometry, microtype, amsmath, amssymb, mathtools, amsthm, graphicx, booktabs, array, url, authblk; the wrapper is the lane's pinned \texttt{amsart} editorial wrapper at 10pt on A4, which restates the theorem-like environments the retained text needs and adds the article's own macro definitions that the retained text needs (none), with extra wrapper packages (bm, bbm, dsfont, xspace, cancel, mathtools). The delivered numbering is the unit's own. Assets: no retained span contains a figure environment, a graphics-inclusion command, a PDF-page inclusion, a table or a TikZ picture; no image file is copied into this package, the package declares no asset dependency and no image file is redistributed with it. Licence: version arXiv:2609.27744v1 of this article is distributed under the Creative Commons Attribution 4.0 International licence, as recorded in the arXiv licence metadata for this exact version and shown on the abstract page https://arxiv.org/abs/2609.27744v1. A full rights notice is delivered at licenses/arxiv-2609.27744-CC-BY-4.0.txt. Characters: every retained excerpt is pure ASCII, so the delivered engine, XeLaTeX with its default Latin Modern text font, reports no missing character.
\begin{equation}
 \Phi_{tk}(s)=\min_{x}\left\{c_{tk}^{\top}x
 +\gamma\widehat Q_{t+1}(s'):\
 s'-s-q_{tk}+u+w=0,\;x\in\mathcal X_{tk}\right\}.
 \label{eq:app-parametric-lp}
\end{equation}

\begin{proposition}[Dual multiplier as a storage subgradient]\label{prop:app-dual}
Assume strong duality holds for \eqref{eq:app-parametric-lp} at a sampled state
$\bar s$. Then
\begin{equation}
 -\pi_{tk}\in\partial\Phi_{tk}(\bar s).
 \label{eq:app-dual-subgrad}
\end{equation}
For a probability-weighted stage value
$Q_t(s)=\sum_kp_k\Phi_{tk}(s)$, a valid subgradient is
\begin{equation}
 \beta_t=-\sum_kp_k\pi_{tk},
 \label{eq:app-beta}
\end{equation}
and the associated Benders cut is
\begin{equation}
 Q_t(s)\ge Q_t(\bar s)+\beta_t(s-\bar s)
 =\alpha_t+\beta_t s.
 \label{eq:app-benders-cut}
\end{equation}
Hence the positive marginal water value under this sign convention is
\begin{equation}
 \lambda_t^W=-\beta_t=\sum_kp_k\pi_{tk}.
 \label{eq:app-discrete-water}
\end{equation}
\end{proposition}

\begin{proof}
The Lagrangian term associated with the storage balance is
$\pi_{tk}(s'-s-q_{tk}+u+w)$. For the optimal multiplier
$\pi_{tk}^*$ at $\bar s$, the dual function evaluated at another parameter $s$
has the affine dependence
\begin{equation}
 d(\pi_{tk}^*;s)=d(\pi_{tk}^*;\bar s)
 -\pi_{tk}^*(s-\bar s).
 \end{equation}
Weak duality gives $\Phi_{tk}(s)\ge d(\pi_{tk}^*;s)$, while strong duality at
$\bar s$ gives $d(\pi_{tk}^*;\bar s)=\Phi_{tk}(\bar s)$. Therefore
\begin{equation}
 \Phi_{tk}(s)\ge\Phi_{tk}(\bar s)
 -\pi_{tk}^*(s-\bar s),
\end{equation}
which proves \eqref{eq:app-dual-subgrad}. Probability-weighted summation yields
\eqref{eq:app-beta}--\eqref{eq:app-benders-cut}. The water value is the reduction
in future cost per additional unit of storage, so it is the negative cut slope,
which gives \eqref{eq:app-discrete-water}. If several cuts are active, the stage
value is nondifferentiable and the admissible marginal values are the negatives of
the active subgradients. This is the discrete analogue of a nonsmooth HJB storage
derivative.
\end{proof}

\end{document}
